\documentclass[12pt, a4paper]{ctexart}

% 严格按照要求重设边距
\usepackage{geometry}
\geometry{left=3.18cm, right=3.18cm, top=2.54cm, bottom=2.54cm}

\usepackage[framemethod=default]{mdframed}
\ctexset{punct=quanjiao}
\xeCJKsetup{CJKecglue = {\hskip 0.1em plus 0.05em minus 0.05em}}
% 仿宋字体定义（专供说明文字使用）
\setCJKfamilyfont{fsgb}[Path=fonts/, AutoFakeBold=1.8]{仿宋_GB2312.ttf}
\newcommand{\fangsonggb}{\CJKfamily{fsgb}}

% 【核心突破】封装说明文字宏，集合仿宋、单倍行距、段后8磅
\newcommand{\instruction}[1]{{\fangsonggb\linespread{1.3}\selectfont\setlength{\parskip}{8bp} #1 \par}}

% 彻底清除所有页码和页眉页脚
\pagestyle{empty}

\begin{document}

% 标题部分
{\noindent\centering\heiti\zihao{4} 预算说明书\par}

\begin{mdframed}[
    linecolor=black,
    linewidth=0.5pt,
    backgroundcolor=white,
    roundcorner=0pt,
    everyline=true,
    userdefinedwidth=16.09cm,
    leftmargin=-0.725cm,
    rightmargin=0.725cm,
    innertopmargin=0cm,
    innerbottommargin=0.2cm,
    innerleftmargin=0.05cm,
    innerrightmargin=0.05cm
]
% 【格式切换】将框内全局基础字体设为宋体（填报的正文将自动应用）
\songti\fontsize{9bp}{9pt}\selectfont
\setlength{\parindent}{2em}

{\bfseries\indent\fangsonggb

1. 科学基金资助项目直接费用\par
请按照《国家自然科学基金项目申请书预算编制说明》等有关要求，按照政策相符性、目标相关性和经济合理性原则，实事求是编制项目/课题预算。填报时，每个科目应结合科研任务按支出用途进行基本测算说明。}

\vspace{0.5em}

\instruction{1.1\hspace{0.75em}设备费（是指在项目/课题实施过程中购置或试制专用仪器设备，对现有仪器设备进行升级改造，以及租赁外单位仪器设备而发生的费用。计算类仪器设备和软件工具可在设备费科目列支。填报时，应按照设备购置费、试制改造费和租赁使用费的分类，提供设备支出的必要性及基本测算说明。单价大于5万元（含5万元）的设备需补充说明设备的主要性能指标、主要技术参数等内容。）}
\vspace{0.3em}

（这里写内容）

\instruction{1.2\hspace{0.75em}业务费（是指项目/课题实施过程中消耗的各种材料、辅助材料等低值易耗品的采购、运输、装卸、整理等费用，发生的测试化验加工、燃料动力、出版/文献/信息传播/知识产权事务、会议/差旅/国际合作交流等费用，以及其他相关支出。填报时，应按照支出大类进行基本测算说明。）}
\vspace{0.3em}

（这里写内容）

\instruction{1.3\hspace{0.75em}劳务费（是指在项目实施过程中支付给参与项目/课题研究的研究生、博士后、访问学者以及项目/课题聘用的研究人员、科研辅助人员等的劳务性费用，以及支付给临时聘请的咨询专家的费用等。填报时，应综合考量劳务费支出对象所承担研究任务的必要性、投入本项目/课题的工作时长、费用标准的合理性等因素，按照人员类别进行基本测算说明。专家咨询费应按照国家有关规定执行。）}
\vspace{0.3em}

（这里写内容）

\vspace{1em}

{\bfseries\indent\fangsonggb 2.直接费用中合作研究转拨资金\par}

\instruction{需对合作研究单位承担的研究任务做必要说明。直接费用转拨资金需经项目/课题申请人与参与者协商一致，并按设备费、业务费、劳务费三个科目做预算说明。如存在多个合作研究单位，请分单位逐一说明。}
\vspace{0.3em}

（这里写内容）

\vspace{1em}

{\bfseries\indent\fangsonggb 3. 其他来源资金\par}

\instruction{对其他来源资金的资金来源、资金具体开支用途做简要说明。}
\vspace{0.3em}

（这里写内容）

\end{mdframed}

\end{document}